%!TEX program = xelatex
\documentclass[border=10pt]{standalone}
\usepackage{fontspec}
\usepackage{polyglossia}
\setmainlanguage{russian}
\setmainfont{FreeSerif}
\usepackage{tikz}
\usetikzlibrary{positioning,shapes.geometric}
\tikzset{
  base/.style={
    rectangle, rounded corners=3pt,
    draw=black!40, thick,
    align=center, font=\small,
    inner xsep=8pt, inner ysep=5pt,
    minimum height=0.75cm,
  },
  result/.style  ={base, fill=black!8,  draw=black!55, font=\small\bfseries},
  proj/.style    ={base, fill=blue!7,   draw=blue!50},
  sel/.style     ={base, fill=teal!10,  draw=teal!55},
  joinnode/.style={base, fill=violet!8, draw=violet!50},
  hashnode/.style={base, fill=gray!10,  draw=gray!50},
  table/.style   ={base, fill=orange!8, draw=orange!50, font=\small\itshape},
  edge/.style    ={->, thick, color=black!50, shorten >=2pt, shorten <=2pt},
  lbl/.style     ={font=\normalsize\bfseries},
}
\begin{document}
\begin{tikzpicture}[node distance=7mm and 14mm]

\node[lbl]    (t2)                           {План №2};
\node[result] (r2)  [below=5mm of t2]        {result};
\node[proj]   (p2)  [below=of r2]
  {$\pi$\ Н\_ЛЮДИ.ОТЧЕСТВО, Н\_ВЕДОМОСТИ.ЧЛВК\_ИД};
\node[joinnode] (j2) [below=of p2]
  {Hash Right Join\\Н\_ЛЮДИ.ИД = Н\_ВЕДОМОСТИ.ЧЛВК\_ИД};

% Left: Seq Scan ВЕДОМОСТИ with filter (probe side)
\node[sel]    (sv) [below left=10mm and 24mm of j2]
  {$\sigma$\ Н\_ВЕДОМОСТИ.ЧЛВК\_ИД $<$ 117219\\
   $\land$\ Н\_ВЕДОМОСТИ.ЧЛВК\_ИД = 117219};
\node[table]  (tv) [below=9mm of sv] {Н\_ВЕДОМОСТИ};

% Right: Hash build on ЛЮДИ
\node[hashnode] (hsh) [below right=10mm and 18mm of j2]
  {Hash\\build};
\node[sel]    (sl)  [below=9mm of hsh]
  {$\sigma$\ Н\_ЛЮДИ.ИД = 100865};
\node[table]  (tl)  [below=9mm of sl] {Н\_ЛЮДИ};

\draw[edge] (p2)  -- (r2);
\draw[edge] (j2)  -- (p2);
\draw[edge] (sv)  -- (j2);
\draw[edge] (hsh) -- (j2);
\draw[edge] (tv)  -- (sv);
\draw[edge] (sl)  -- (hsh);
\draw[edge] (tl)  -- (sl);

\node[above=3mm of t2, font=\footnotesize, text=black!50]
  {Фильтры применяются до соединения; хеш строится по Н\_ЛЮДИ (1 строка)};
\end{tikzpicture}
\end{document}